% ============================================================
\chapter{Conclusions}
\label{ch:conclusions}
% ============================================================

This work set out to determine how far monocular RGB segmentation can support
vehicle-specific traversability perception under realistic computational
constraints. CaT provided a suitable test because its labels encode nested
vehicle capability rather than terrain material alone. Auditing the release
identified 1,812 valid RGB--label pairs. Their integer maps were converted into
explicit blue-only and blue+green policies, and the official 544-image test
partition was kept separate from the deterministic 1,002/266 development
split.

PIDNet-S provided the first complete answer to the implementation question. The
selected 7,623,522-parameter model achieves 0.893941 mIoU, with 4.32\% FSR and
7.04\% FBR. Reconstructing its 133,693,440-pixel confusion matrix on the CPU
confirms that the saved configuration and checkpoint reproduce the reported
result. Its development also produced negative findings of practical value:
the evaluated augmentation, false-safe penalty, and higher-resolution setting
did not improve the reference configuration. The threshold sweep showed that
conservativeness can be adjusted, but only by exchanging one disagreement
direction for the other.

Those disagreements acquire meaning only when their spatial form is inspected.
The large false-safe field, the almost removed shadowed corridor, and the
shifted vegetation boundary would lead to different planning responses despite
all appearing as pixel errors. Percentile sampling shows that such behaviour
is not represented adequately by a few best and worst images. ALICE sequences
add evidence about output behaviour under domain shift, but without labels
they cannot supply an accuracy estimate.

The comparative study demonstrates both the value and the limits of pretrained
transfer. In the original fixed-15-epoch matrix, ROD raises the matched
blue+green mean from PIDNet-S's 0.8822 to 0.9138. The convergence-aware
follow-up shows that neither value represented the end of learning: they rise
to 0.9195 and 0.9331. Every architecture improves beyond its 15-epoch result,
and validation selects checkpoints between epochs 50 and 141. The largest
mean is $0.9423\pm0.0005$ for FPN/EfficientNet-B0, which ranks first in all
three executed seeds; U-Net/EfficientNet-B0 follows at
$0.9409\pm0.0004$. This identifies the highest observed mean for the declared suite
without claiming that one shared schedule is optimal for every architecture.

No single ordering survives the change from accuracy to runtime. PIDNet-S is
fastest on the common H100 protocol at 270.38 FPS and reaches 106.43 FPS in the
RTX~5060 TensorRT pipeline. DDRNet-23-Slim leads the compiled Ryzen measurement
at 36.50 FPS, while the most accurate EfficientNet-B0 models are slower and ROD
reaches only 0.46 compiled CPU FPS.
Because training changes parameter values rather than the inference graph,
these architecture-level timings remain applicable to the convergence-aware
accuracy points.
The TensorRT-compiled ROD model reaches 39.38 FPS with negligible mIoU change. These measurements
show why hardware, precision, runtime, and timing boundaries belong in the
model specification.

The four research questions of Chapter~\ref{ch:introduction} can now be
answered directly:
\begin{enumerate}
  \item PIDNet-S does provide a reproducible, real-time baseline. Trained only
  on the official training partition, the retained checkpoint reaches
  0.893941 mIoU with 4.32\% FSR and 7.04\% FBR, its complete confusion matrix
  is reproduced exactly over 133,693,440 pixel decisions, and the same saved
  parameters sustain 21.16 eager CPU frames per second and a 106.43 FPS
  TensorRT pipeline.
  \item The tested choices that materially change accuracy and error balance
  are augmentation (disabling the evaluated recipe improves mIoU by 2.72
  points), class weighting (neutral weights lower FBR to 6.95\% but raise
  FSR), and the decision threshold (it exchanges FSR for FBR monotonically).
  The false-safe loss penalty and the higher-resolution variant degrade the
  operating point.
  \item Under matched seeds, frozen-encoder transfer outperforms the compact
  trained-from-scratch network: ROD leads PIDNet-S by 3.15 percentage points
  for blue+green and 3.79 for blue-only, with the higher score in all six
  policy--seed runs. Within the nine-system matrix, both rank below the
  ImageNet-pretrained encoder--decoder systems, and FPN/EfficientNet-B0 leads
  the convergence-aware suite at $0.9423\pm0.0005$.
  \item Accuracy rankings change substantially once deployment conditions
  enter the comparison. PIDNet-S is fastest under the common H100 protocol at
  270.38 FPS and reaches 106.43 FPS in the TensorRT pipeline; compiled
  DDRNet-23-Slim leads the Ryzen measurement at 36.50 FPS; and the most
  accurate EfficientNet-B0 systems sit in the slower half of the table,
  confirming that parameter counts and architecture families predict measured
  speed poorly.
\end{enumerate}

The resulting masks are useful perception outputs, not certificates of safe
motion. A single RGB image cannot reveal every property that determines vehicle
support, and the available CaT release contains neither repeated annotations
nor agreement measurements for ambiguous capability boundaries. The study also
does not test temporal consistency, planning, vehicle dynamics, or the complete
embedded stack. Progress beyond the present results therefore depends less on
adding another architecture row than on collecting stronger target evidence,
labelling cross-domain sequences, modelling uncertainty, and evaluating the
perception--planning--vehicle loop as a whole.
